\documentclass[11pt,a4paper]{article}

\usepackage[utf8]{inputenc}
\usepackage[T1]{fontenc}
\usepackage[french]{babel}
\usepackage{geometry}
\geometry{margin=2.5cm}
\usepackage{graphicx}
\usepackage{booktabs}
\usepackage{hyperref}

\title{Cross-Domain Visual Anomaly Detection\\
\large One pipeline, three domains: industry, healthcare and aerial imagery}
\author{Hakim}
\date{\today}

\begin{document}
\maketitle

\begin{abstract}
Ce projet construit un pipeline unique de détection d'anomalies visuelles
\emph{normal-only} appliqué à trois domaines (industrie, santé, aérien). Le
système apprend uniquement sur des images normales, produit un score
d'anomalie, une décision seuillée et une carte de chaleur.
\textbf{Résultats : À compléter après expérimentation.}
\end{abstract}

\section{Problème et motivation}
Les anomalies sont rares et coûteuses à annoter. Apprendre le
« normal » puis mesurer l'écart est souvent plus réaliste que collecter tous
les défauts possibles. \emph{À compléter.}

\section{Données}
MVTec AD (industrie), RSNA Pneumonia (santé, image-level), xView2/xBD
(aérien). Sources, licences et méthodes de séparation : \emph{À compléter.}

\section{Méthode}
Preprocessing, autoencodeur convolutif, PatchCore, score d'anomalie,
génération de heatmap, sélection du seuil. Schéma du pipeline : \emph{À
compléter.}

\section{Résultats}

\subsection{Protocole et métriques}
Le seuil de décision est choisi sur la \textbf{validation} (percentile des
scores d'images normales), jamais sur le test. Les stratégies F1 et coût sont
calibrées sur une part disjointe du test. Métriques niveau image : AUROC,
AUPRC, précision, rappel, F1, matrice de confusion, taux de faux positifs
(FPR) et faux négatifs (FNR). Niveau pixel (si masque fiable) : AUROC/AUPRC
pixel, Dice, IoU.

\subsection{Métriques image-level}
\begin{table}[h]
\centering
\begin{tabular}{lccccc}
\toprule
Modèle & AUROC & AUPRC & Précision & Rappel & F1 \\
\midrule
Autoencodeur & \emph{à compl.} & \emph{à compl.} & \emph{à compl.} & \emph{à compl.} & \emph{à compl.} \\
PatchCore    & \emph{à compl.} & \emph{à compl.} & \emph{à compl.} & \emph{à compl.} & \emph{à compl.} \\
\bottomrule
\end{tabular}
\caption{Résultats image-level (MVTec AD). \emph{À compléter après expérimentation.}}
\end{table}

\subsection{Métriques pixel-level}
\begin{table}[h]
\centering
\begin{tabular}{lcccc}
\toprule
Modèle & AUROC pixel & AUPRC pixel & Dice & IoU \\
\midrule
Autoencodeur & \emph{à compl.} & \emph{à compl.} & \emph{à compl.} & \emph{à compl.} \\
PatchCore    & \emph{à compl.} & \emph{à compl.} & \emph{à compl.} & \emph{à compl.} \\
\bottomrule
\end{tabular}
\caption{Résultats pixel-level (masques disponibles). \emph{À compléter.}}
\end{table}

\subsection{Comparaison des modèles (qualité et coût)}
\begin{table}[h]
\centering
\begin{tabular}{lccc}
\toprule
Critère & Autoencodeur & PatchCore \\
\midrule
AUROC image & \emph{à compl.} & \emph{à compl.} \\
Temps d'inférence (ms/img) & \emph{à compl.} & \emph{à compl.} \\
Empreinte mémoire (banque, Mo) & 0 & \emph{à compl.} \\
Entraînement & gradient (époques) & construction de mémoire \\
Qualité des heatmaps & \emph{à compl.} & \emph{à compl.} \\
\bottomrule
\end{tabular}
\caption{Compromis qualité / coût. Généré par \texttt{scripts/compare\_models.py}.}
\end{table}

Analyse des erreurs (faux positifs, faux négatifs) et cartes de chaleur :
figures générées dans \texttt{reports/figures/}. \emph{À compléter.}

\section{Transférabilité}

\textbf{Ce qui est réellement commun} entre les trois domaines : le
preprocessing, l'autoencodeur, l'entraînement, l'évaluation et l'inférence sont
strictement identiques (mêmes classes, même code). \textbf{Ce qui est adapté}
tient dans la configuration YAML et de fines classes dataset : la lecture des
données et surtout la \emph{règle de découpage}, propre à chaque domaine.

\begin{table}[h]
\centering
\begin{tabular}{lccc}
\toprule
Domaine & Règle de split & Masques pixel & AUROC \\
\midrule
Industrie (MVTec) & aléatoire (normales) & oui & \emph{à compl.} \\
Santé (RSNA) & \textbf{par patient} & non & \emph{à compl.} \\
Aérien (xView2) & \textbf{par zone} & oui & \emph{à compl.} \\
\bottomrule
\end{tabular}
\caption{Comparaison inter-domaines. Généré par \texttt{scripts/run\_cross\_domain.py}.
\emph{À compléter après expérimentation.}}
\end{table}

Facteurs d'impact à analyser : textures (industrie très régulière vs aérien
très variable), résolution, type d'anomalie (défaut localisé vs changement
diffus) et variabilité du « normal ». La séparation \textbf{par patient}
(santé) et \textbf{par zone} (aérien) est indispensable pour éviter des scores
faussement optimistes. \emph{À compléter.}

\section{Architecture et MLOps}
Données $\rightarrow$ validation $\rightarrow$ preprocessing $\rightarrow$
entraînement $\rightarrow$ suivi MLflow $\rightarrow$ évaluation $\rightarrow$
export $\rightarrow$ API $\rightarrow$ interface $\rightarrow$ monitoring.

\section{Limites et éthique}
Biais des datasets, changement de domaine, faux négatifs, interprétation
excessive des heatmaps, confidentialité, RGPD, anonymisation. \textbf{Le
module médical n'est pas un dispositif de diagnostic.}

\section{Compétences démontrées}
Computer Vision, PyTorch, détection d'anomalies, évaluation,
interprétabilité, API, Docker, CI, tests, MLOps, communication scientifique.

\section{Questions d'entretien}
\emph{À compléter (8 à 10 questions/réponses).}

\end{document}
